% Source: https://artofproblemsolving.com/community/c3h1968263_help_2019_amc_10a_10
% Stable identity: aops-topic-1968263-q01
Put the floor on coordinates with corners $(0,0)$ and $(17,10)$; the bug walks along $y=\dfrac{10}{17}x$.

\textbf{The path never passes through an interior grid point.} If it passed through a lattice point $(x,y)$ with $0<x<17$, then $17y=10x$; since $\gcd(10,17)=1$, $17$ would divide $x$, which is impossible. So the bug never crosses a vertical and a horizontal line at the same moment.

\textbf{Counting tiles.} The bug starts inside the first tile. Every time it crosses one interior grid line it leaves the tile it is in and enters a new one, and it enters a new tile at no other time. So
\[
\text{tiles visited}=1+(\text{number of crossings}).
\]
The ``$+1$'' is the starting tile, which the bug is in before it has crossed anything (just as a fence of $k$ gaps has $k+1$ posts). The bug crosses each of the $16$ interior vertical lines $x=1,\dots,16$ and each of the $9$ interior horizontal lines $y=1,\dots,9$ exactly once, and these $25$ crossings happen at different moments. Hence the bug visits
\[
1+16+9=\boxed{26}
\]
tiles (shaded below).

\begin{asy}
usepackage("amsmath");
size(7cm);
defaultpen(fontsize(10pt));
fill(box((0,0),(1,1)),palegray);
fill(box((1,0),(2,1)),palegray);
fill(box((1,1),(2,2)),palegray);
fill(box((2,1),(3,2)),palegray);
fill(box((3,1),(4,2)),palegray);
fill(box((3,2),(4,3)),palegray);
fill(box((4,2),(5,3)),palegray);
fill(box((5,2),(6,3)),palegray);
fill(box((5,3),(6,4)),palegray);
fill(box((6,3),(7,4)),palegray);
fill(box((6,4),(7,5)),palegray);
fill(box((7,4),(8,5)),palegray);
fill(box((8,4),(9,5)),palegray);
fill(box((8,5),(9,6)),palegray);
fill(box((9,5),(10,6)),palegray);
fill(box((10,5),(11,6)),palegray);
fill(box((10,6),(11,7)),palegray);
fill(box((11,6),(12,7)),palegray);
fill(box((11,7),(12,8)),palegray);
fill(box((12,7),(13,8)),palegray);
fill(box((13,7),(14,8)),palegray);
fill(box((13,8),(14,9)),palegray);
fill(box((14,8),(15,9)),palegray);
fill(box((15,8),(16,9)),palegray);
fill(box((15,9),(16,10)),palegray);
fill(box((16,9),(17,10)),palegray);
for(int i=0;i<=17;++i) draw((i,0)--(i,10),gray(0.55));
for(int j=0;j<=10;++j) draw((0,j)--(17,j),gray(0.55));
draw((0,0)--(17,0)--(17,10)--(0,10)--cycle,linewidth(1));
draw((0,0)--(17,10),red+linewidth(1.2));
label("$17$",(8.5,0),S);label("$10$",(17,5),E);
label("\mbox{}",(min(currentpicture,user=true)+max(currentpicture,user=true))/2,black); // keeps the text after the figure black when LaTeX typesets the labels inline
\end{asy}

(In general an $m\times n$ grid diagonal meets $m+n-\gcd(m,n)$ cells.)
